\begin{figure}[H]
	\centering
	\begin{tikzpicture}[node distance=2cm]

	% 节点
	% \node[startstop] (start) {开始};
	% \node[block, below=of start] (step1) {步骤 1};
	% \node[decision, below=of step1] (decide) {条件？};
	% \node[block, below left=2cm and 1.5cm of decide] (step2a) {路径 A};
	% \node[block, below right=2cm and 1.5cm of decide] (step2b) {路径 B};
	% \node[startstop, below=3cm of decide] (end) {结束};
	\node[block] (车速传感器) {车速传感器};
	\node[startstop,below=of 车速传感器] (控制器) {控制器};
	\node[block,right=of 控制器] (排肥阻力传感器) {排肥阻力传感器};
	\node[block,above=of 排肥阻力传感器] (肥箱肥量传感器) {肥箱肥量传感器};
	\node[block,below=of 排肥阻力传感器] (流量传感器) {流量传感器};
	\node[startstop,right=of 排肥阻力传感器] (排肥装置) {排肥装置};
	\node[block,below=of 排肥装置] (超声波换能器) {超声波换能器};

	% 连接
	% \draw[arrow] (start) -- (step1);
	% \draw[arrow] (step1) -- (decide);
	% \draw[arrow] (decide) -- node[left] {是} (step2a);
	% \draw[arrow] (decide) -- node[right] {否} (step2b);
	\draw[arrow] (车速传感器) -- node[fill=white]{前进速度} (控制器);
	\draw[arrow] (肥箱肥量传感器) -- node[fill=white]{肥量} (控制器);
	\draw[arrow] (排肥阻力传感器) -- node[fill=white]{力} (控制器);
	\draw[arrow] (流量传感器) -- node[fill=white]{排肥速度} (控制器);
	\draw[arrow] (控制器) |- node[fill=white]{功率} (超声波换能器);
	\draw[arrow] (超声波换能器) -- (排肥装置);
	\draw[arrow] (排肥装置) -- (肥箱肥量传感器);
	\draw[arrow] (排肥装置) -- (排肥阻力传感器);
	\draw[arrow] (排肥装置) -- (流量传感器);

	% 两路径指向结束
	% \draw[arrow] (step2a) |- (end);
	% \draw[arrow] (step2b) |- (end);

	\end{tikzpicture}
	\caption{控制系统示意图}
	\label{pic:控制系统示意图}
\end{figure}